% ID: FIS-TER-GAS-002
% Titolo: Compressione isobara ed espansione isoterma di un gas
% Disciplina: Fisica
% Area: Termologia
% Argomento: Gas perfetti
% Sottoargomento: Trasformazioni isobare e isoterme
% Classe: Quarta liceo scientifico
% Difficolta: 4
% Tipo: problema_numerico
% Risultato: T_B = 187 K; V_C = 2,60 L; trasformazione non ciclica
% Tempo_stimato: 12 min
% Competenze: modellizzazione, uso dell'equazione di stato, lettura del piano pressione-volume
% Prerequisiti: gas perfetti, trasformazione isobara, trasformazione isoterma
% Tag: termologia, gas_perfetti, isobara, isoterma, piano_pv
% Fonte: verifica_fisica_gas_perfetti_anonimizzata
% Autore: Riccardo Carli
% Licenza: CC-BY-SA-4.0

\begin{esercizio}
Un gas perfetto occupa un volume iniziale di
\(2{,}00 \times 10^{-3}\,\mathrm{m^3}\) alla temperatura di
\(15\,^\circ\mathrm{C}\), nello stato \(A\). Il gas viene compresso di
\(0{,}70\,\mathrm{L}\), mantenendo la pressione costante, fino allo stato
\(B\). Dopo la compressione il gas viene mantenuto a temperatura costante e la
pressione viene dimezzata, raggiungendo lo stato \(C\).

\begin{enumerate}
  \item Disegna il grafico pressione-volume delle trasformazioni subite dal gas.
  \item Stabilisci se la trasformazione complessiva e ciclica.
  \item Calcola la temperatura del gas dopo la compressione, nello stato \(B\).
  \item Calcola il volume finale occupato dal gas, nello stato \(C\).
\end{enumerate}
\end{esercizio}

\begin{soluzione}
Il volume iniziale e:
\[
V_A=2{,}00\cdot 10^{-3}\,\mathrm{m^3}=2{,}00\,\mathrm{L}.
\]
La compressione isobara riduce il volume di \(0{,}70\,\mathrm{L}\), quindi:
\[
V_B=2{,}00\,\mathrm{L}-0{,}70\,\mathrm{L}
=1{,}30\,\mathrm{L}.
\]

Nel grafico pressione-volume il tratto \(A\to B\) e un segmento orizzontale
verso sinistra, perche la pressione resta costante e il volume diminuisce. Il
tratto \(B\to C\) e un'isoterma: poiche la pressione si dimezza, il volume
aumenta.

La trasformazione complessiva non e ciclica, perche lo stato finale \(C\) non
coincide con lo stato iniziale \(A\).

Durante la trasformazione isobara vale:
\[
\frac{V_A}{T_A}=\frac{V_B}{T_B}.
\]
La temperatura iniziale in kelvin e:
\[
T_A=15+273=288\,\mathrm{K}.
\]
Quindi:
\[
T_B=T_A\frac{V_B}{V_A}
=288\,\mathrm{K}\cdot \frac{1{,}30}{2{,}00}
\simeq 187\,\mathrm{K}.
\]

Nel tratto isoterma \(B\to C\), \(pV=\text{costante}\). Se la pressione viene
dimezzata, il volume raddoppia:
\[
V_C=2V_B=2{,}60\,\mathrm{L}
=2{,}60\cdot 10^{-3}\,\mathrm{m^3}.
\]
\end{soluzione}
